% !PRINTED 600
\weberpage{600}

Unter der Voraussetzung also, dass der Satz IV für Gleichungen $\mu^{\text{ten}}$ Grades richtig ist, folgt seine Richtigkeit für Gleichungen $\mu m^{\text{ten}}$ Grades; und da er für Gleichungen von Primzahlgrad gilt, so ist er allgemein nachgewiesen.

\begin{center}
\S.\,177.

\gesperrt{Einfachheit der alternirenden Gruppe.}
\end{center}

Wir haben früher gesehen (\S.\,149), dass, wenn wir die Co\"efficienten einer Gleichung $n^{\text{ten}}$ Grades als unabhängige Variable und den Körper aller rationalen Functionen dieser Co\"efficienten als Rationalitätsbereich betrachten, die \gesperrt{Galois}'sche Gruppe der Gleichung die symmetrische Gruppe ist. In der symmetrischen Gruppe ist immer ein Normaltheiler vom Index 2 enthalten, die \gesperrt{alternirende Gruppe}, auf die sich die Gruppe der Gleichung reducirt, wenn die Quadratwurzel aus der Discriminante adjungirt wird. Bei vier Ziffern hat die alternirende Gruppe, die aus der identischen Permutation, acht dreigliedrigen Cykeln und drei Paaren von Transpositionen besteht, den Normaltheiler vom Index 3:

\[
1, \quad (0,\,1)(2,\,3), \quad (0,\,2)(1,\,3), \quad (0,\,3)(1,\,2).
\]

Diese Gruppe hat drei verschiedene Normaltheiler vom Index 2, von denen wir einen 1, $(0,\,1)(2,\,3)$ bevorzugen, von dem wieder die Einheitsgruppe ein Normaltheiler vom Index 2 ist. Die Gruppe der 24 Permutationen von vier Ziffern ist also metacyklisch, und darauf beruht jede Auflösungsmethode der biquadratischen Gleichung (\S.\,160, 161).

Wir wollen nun nachweisen, dass, wenn $n$ grösser als 4 ist, die alternirende Gruppe ausser der Einheitsgruppe überhaupt keine normalen Theiler hat, oder nach der früher eingeführten Bezeichnung \gesperrt{einfach} ist. Daraus folgt dann, dass die Bedingung, die wir für die algebraische Auflösbarkeit einer Gleichung als nothwendig gefunden haben, für die Gleichungen von höherem als dem vierten Grade, deren Gruppe die symmetrische oder die alternirende ist, nicht erfüllt ist, und dass also \gesperrt{Gleichungen von höherem als dem vierten Grade, so lange die Co\"efficienten unabhängige Variable sind, nicht mehr algebraisch lösbar sind.}
